\chapter{Geltungsbereich und Grundbegriffe}
\label{b08guide:ch:foundations}

Dieser Leitfaden ergänzt Band~8 um zentrale Definitionen und Sätze
über bijektive Funktionen zwischen beliebigen Mengen. Grundlage
sind ZF ohne Auswahlaxiom und der Funktionsbegriff aus Band~5.
Beweise und Beweisskizzen entfallen.

\section{Gegenstand und Abgrenzung}

Im Mittelpunkt stehen Charakterisierungen der Bijektivität,
Umkehrbarkeit und Konstruktionen aus gegebenen Bijektionen.
Allgemeine Funktionssätze und Resultate, für die schon
Injektivität oder schon Surjektivität genügt, bleiben in den
Bänden~5--7. Äquivalenz- und Ordnungsrelationen, Mächtigkeitsrechnung,
endliche Abzählformeln und zusätzliche algebraische oder topologische
Strukturen werden nicht behandelt. Auch Cantor--Schröder--Bernstein
und die speziellen Potenzmengenkonstruktionen des Hauptbands entfallen.

Eine Selbstbijektion wird nur als Funktion auf einer Menge betrachtet;
Mengeninklusionen und logische Gleichwertigkeiten dienen allein der
Formulierung. Bei Familien von Bijektionen ist die ganze Familie
gegeben. Es wird keine gleichzeitige Auswahl aus Existenzaussagen
vorausgesetzt. Leere Mengen sind zugelassen.

Die Zusammenstellung wählt bekannte Hauptresultate aus und
beansprucht keine vollständige Erfassung aller bekannten Sätze.
Die Gestaltung folgt den Leitfäden zu Band~3, Band~4, Band~5,
Band~28 und Band~37.

\section{Bijektivität und Notation}

\begin{definition}[Bijektive Funktion und Bijektionsmenge]
\label{b08guide:def:bijection}
Eine Funktion $f\colon A\to B$ heißt \emph{bijektiv von $A$ auf $B$},
wenn jedes Element von $B$ genau ein Urbild in $A$ besitzt:
\[
 \forall b\in B\;\exists!\,a\in A\quad f(a)=b.
\]
Wir schreiben $f\colon A\bij B$. Mit
$B^A=\{h\mid h\colon A\to B\}$ sei
\[
 \operatorname{Bij}(A,B)
 =\{f\in B^A\mid f\colon A\bij B\}.
\]
\end{definition}

Eine Funktion ist wie im Hauptband ihr Graph. Die Aussage der
Bijektivität bezieht sich stets auf die angegebenen Mengen $A,B$.

\begin{definition}[Bild, Urbild und abgeleitete Abbildungen]
\label{b08guide:def:notation}
Für $f\colon A\to B$, $U\subseteq A$ und $V\subseteq B$ sei
\[
 \begin{aligned}
 f[U]&=\{f(a)\mid a\in U\},&
 f^{-1}[V]&=\{a\in A\mid f(a)\in V\},\\
 f_*\colon\mathcal P(A)&\to\mathcal P(B),\quad U\mapsto f[U],&
 f^*\colon\mathcal P(B)&\to\mathcal P(A),\quad V\mapsto f^{-1}[V].
 \end{aligned}
\]
Hier bezeichnet $\mathcal P(A)$ die Potenzmenge von $A$.
Die Schreibweise $f^{-1}[V]$ ist für jede Funktion eine Urbildmenge;
sie setzt noch keine Umkehrfunktion voraus.
\end{definition}

Als Hilfsnotation aus Band~5 verwenden wir $f|_U$ für die
Einschränkung, $(g\circ f)(a)=g(f(a))$ für die Komposition
($f\colon A\to B$, $g\colon B\to C$) und
$\operatorname{id}_A(a)=a$ für die Identität. Funktionswerte
werden nur im jeweiligen Definitionsbereich verwendet.

\section{Wegweiser durch die Resultate}

\begin{center}
\begin{tabularx}{\textwidth}{@{}p{.31\textwidth}X@{}}
\toprule
Fragestellung & Fundstelle\\
\midrule
Bijektivität erkennen &
Umkehrbarkeit, Teilmengen- und Komplementkriterien
in Kapitel~\ref{b08guide:ch:criteria}.\\
Mit Umkehrungen rechnen &
Umkehrgesetze, Komposition, Zwei-aus-drei-Regel und
Funktionsgleichungen in Kapitel~\ref{b08guide:ch:inverses}.\\
Bereiche verändern &
Vollständige Teilmengenkorrespondenz, Einschränkungen,
Verkleben und Erweiterungen in Kapitel~\ref{b08guide:ch:restrictions}.\\
Neue Bijektionen bilden &
Markierte Vereinigungen, Produkte und Umindizierung
in Kapitel~\ref{b08guide:ch:constructions}.\\
Funktionen umbenennen &
Wechsel von Argumenten und Werten, Bijektionsmengen
und ein Zweipunkttest in Kapitel~\ref{b08guide:ch:functionsets}.\\
\bottomrule
\end{tabularx}
\end{center}

Die Nummern mit Präfix $8H$ gehören nur zu diesem Leitfaden.
Die Nummerierung des vollständigen Fachbands bleibt unverändert.
